\section{Boolean functions implemented by real-valued tanh teachers}
\label{sec:analog-teacher-critical-rate}

The circuit-rate bound also applies to Boolean targets specified by a
bounded-weight tanh teacher. Finite precision is enough because the
teacher's final output has a fixed margin. The necessary precision
grows linearly with teacher depth, and each numerical layer has a
constant-depth threshold implementation of polynomial size.
No stability condition on the teacher's hidden states is needed.

\begin{lemma}[Finite-precision threshold compilation]
\label{lem:analog-teacher-threshold-compilation}
There is a universal integer $a$ with the following property.
Suppose a tanh network on $\{-1,1\}^d$ has depth $D\ge1$, hidden
width at most $C$, parameter magnitudes at most $B\ge1$, and a linear
readout $z_T$ satisfying
\[
 f(x)z_T(x)\ge\tfrac12\qquad\hbox{for every }x,
 \qquad f:\{-1,1\}^d\to\{-1,1\}.
\]
Depth counts affine maps, including the linear readout. Biases may be
present; the original model with a constant input and no internal
biases is a special case.
Then $f$ has a layered Boolean threshold implementation of depth
$h\le a(D+1)$, width at most $K^a$, and gate coefficient Euclidean
norm at most $K^a$, where
\[
 K=d+C+B+D+2.
\]
Every Boolean gate has margin at least one, and constant rails and
layering overhead are included in these bounds.
\end{lemma}

\begin{proof}
We use the standard arithmetic threshold-circuit theorem: addition,
comparison, division, and iterated multiplication of polynomially many
polynomial-bit integers have constant-depth, polynomial-size circuits
of unweighted majority and Boolean gates; see
Hesse--Allender--Barrington~\cite{hesse2002} for the nontrivial division
and iterated-multiplication assertions. The depth and polynomial degree
are universal. Thus this theorem provides majority gates, rather than
arbitrary threshold gates with uncontrolled coefficient magnitudes.

First consider approximating a single scalar tanh.
For a dyadic input $u$ with $|u|\le R$ and desired error $2^{-p}$,
put $t=2|u|$ and
\[
 E_m(t)=\sum_{j=0}^{m-1}\frac{t^j}{j!},
 \qquad m=\lceil40(R+p+2)\rceil.
\]
The Taylor remainder satisfies
\[
 0\le e^t-E_m(t)
 \le e^{2R}\frac{(2R)^m}{m!}
 \le e^{2R}\left(\frac{2eR}{m}\right)^m
 <2^{-p-4}.
\]
Write the sum over a common denominator. All its powers, factorial
constants, integer products, and sums involve polynomially many bits
in $R$, $p$, and the input bit length. They can therefore be computed
in constant threshold depth by the arithmetic theorem. Since
$E_m(t)\ge1$, the map $v\mapsto(v-1)/(v+1)$ has derivative at most
one for $v\ge1$. Division followed by rounding to $p+2$ fractional
bits gives
\[
 \operatorname{sign}(u)\frac{E_m(2|u|)-1}{E_m(2|u|)+1}
\]
to total error less than $2^{-p}$ from $\tanh u$.
Choose the rounding convention to keep the result in $[-1,1]$.
This is a constant-depth majority circuit whose size is polynomial
in the stated numerical parameters, with universal degree.

For the teacher, set
\[
 K_0=4(d+C+1)(B+1),\qquad
 p=\left\lceil(D+3)\log_2 K_0\right\rceil+5.
\]
Round every parameter to $p$ fractional bits and hardcode these bits
in the circuit. This is a nonuniform existence statement and requires
no computability assumption on the original real parameters.
Starting from the exact
Boolean input, evaluate each affine map on the stored dyadic values
exactly, and evaluate each tanh to error at most $2^{-p}$ by the
preceding construction. All stored hidden values have magnitude at
most one, so every preactivation has magnitude at most
$R=(d+C+1)(B+1)$.
Parameter rounding contributes at most $(d+C+1)2^{-p}$ at any layer.
Since tanh is one-Lipschitz, a conservative induction bounds the
error after $j$ affine maps by $K_0^{j+1}2^{-p}$, including the linear
readout. The final error is therefore less than $1/16$.
Its sign is exactly $f$ by the margin assumption.

The arithmetic at every layer has constant threshold depth and
polynomial size in $d,C,B,D$: the precision $p$, the range $R$, all
intermediate bit lengths, and the number of Taylor terms are polynomial
in these parameters. Stack these blocks over $D$ layers and append a
sign comparison. The resulting depth is at most $a(D+1)$ and its
total size is at most $K^a$ after increasing a universal $a$.
Make it layered by identity gates and propagate a constant wire.
This introduces only polynomial overhead and a universal constant
factor in depth.

Finally, an unweighted majority or Boolean gate with fan-in at most
the total circuit size can be written as an integer threshold gate
with coefficients of polynomial magnitude and margin at least one.
For example, place the threshold halfway between the two neighboring
allowed integer sums, and multiply the score by two if necessary.
The constant rail carries the threshold. Its coefficient norm is
polynomial in the circuit size. Increasing $a$ once more proves all
three bounds simultaneously.
\end{proof}

\begin{corollary}[Polynomial teacher-depth dependence of the critical rate]
\label{cor:analog-teacher-critical-rate}
For a Boolean target with the tanh-teacher implementation in
Lemma~\ref{lem:analog-teacher-threshold-compilation}, there is a
universal $A$ such that, whenever
\[
 L\ge A(D+1)\log(d+C+B+D+L+2),
\]
the original unrestricted Gaussian learner satisfies
\[
 J_L(f)\le V_*+
 A\frac{(D+1)^2\log^2(d+C+B+D+L+2)}L.
\]
In particular, every
\[
 L\ge A(D+1)^2\log^2(d+C+B+D+2)
\]
puts $f$ below a fixed universal absolute critical-rate budget.
\end{corollary}

\begin{proof}
Apply the compiler lemma. Its threshold depth is $h\le a(D+1)$,
while its state width and normalized gate coefficient norms are
polynomial in $K=d+C+B+D+2$.
The parameter $T$ in Theorem~\ref{thm:controlled-circuit-rate}
is consequently at most $A_0(D+1)\log(K+L)$.
Increase the universal $A$ so that the displayed depth condition
ensures $L-h\ge L/2$ and $L-h\ge C_0T$.
The theorem gives the rate estimate. Its fixed-budget conclusion,
or the same squared-logarithm comparison used in its proof, gives
the final assertion.
\end{proof}

The result concerns the rate of the target's full Boolean sign pattern
under the original Gaussian prior. It removes an exponential
teacher-depth cost from this rate upper bound. Corollary~\ref{cor:critical-analog-teacher-learning} combines it with
the capacity theorem to give a sample bound polynomial in learner depth
for each fixed teacher depth. Joint polynomial dependence on teacher
width and depth remains open.
